%% kapa: exam=1 round=26 year=2015 period=1 form=A subject=economics subject_ko=경제원론 seq=11 number=51
%% source: Q-Net 감정평가사 기출문제 (cst00309) artlSeq=5138001 file=extracted/1008/5138001/제26회_감정평가사 1차_기출/제26회_감정평가사 1차_기출/1교시_감평_A형.pdf
%% conversion: kordoc 4.24.51 -> tools/gen1.py (auto); review=proofread-sonnet; flags=none
\begin{question}{51}
소득-여가 결정모형에서 효용극대화를 추구하는 甲의 노동공급에 관한 설명으로 옳은 것은? (단, 소득과 여가는 모두 정상재이며, 소득효과 및 대체효과의 크기 비교는 절대값을 기준으로 한다.)
\begin{choices}
  \item 시간당 임금이 상승할 경우, 대체효과는 노동공급 감소요인이다.
  \item 시간당 임금이 상승할 경우, 소득효과는 노동공급 증가요인이다.
  \item 시간당 임금이 하락할 경우, 소득효과와 대체효과가 동일하다면 노동공급은 감소한다.
  \item 시간당 임금이 하락할 경우, 소득효과가 대체효과 보다 크다면 노동공급은 증가한다.
  \item 시간당 임금의 상승과 하락에 무관하게 소득과 여가가 결정된다.
\end{choices}
\end{question}
\answer{4}
